% ========================================================
% RECURSOS DE CÓMPUTO Y PUNTO DE EQUILIBRIO
% ========================================================
\section{Recursos de cómputo y punto de equilibrio}\label{sec:recursos}

\subsection{Memoria y CPU por configuración}

La Tabla~\ref{tab:recursos} recoge memoria y CPU de cada configuración en la laptop del equipo,
sobre el dataset completo, con \texttt{tp/scripts/medir\_recursos.sh}. Acá alcanza una corrida por
configuración: el consumo es estable, y el tiempo, que sí varía, es lo que el benchmark mide con
repeticiones.

% Generado por tp/scripts/tablas_informe.py desde tp/reports/*.json. No editar a mano.
\begin{table}[H]
\caption{Uso de recursos por configuración en la laptop}\label{tab:recursos}
\footnotesize
\begin{tabular}{@{}lrrrrrr@{}}
\toprule
\textbf{Modo} & \textbf{$P$} & \textbf{Clust. (ms)} & \textbf{Heap (MB)} & \textbf{RSS (MB)} & \textbf{CPU (s)} & \textbf{Núcleos} \\
\midrule
seq & --- & 2.344 & 349 & 453 & 5,05 & 1,2 \\
conc & 1 & 2.708 & 349 & 466 & 6,11 & 1,2 \\
conc & 2 & 1.451 & 349 & 452 & 6,26 & 1,7 \\
conc & 4 & 813 & 349 & 532 & 6,71 & 2,1 \\
conc & 8 & 742 & 349 & 531 & 9,07 & 3,0 \\
conc & 16 & 741 & 349 & 531 & 9,07 & 3,0 \\
\bottomrule
\end{tabular}

{\footnotesize\textit{Nota.} fedora · go1.26.7-X:nodwarf5 · 8 CPUs lógicas (4 núcleos físicos con SMT) · dataset completo, $k=8$, 10 iteraciones. Una corrida por configuración: acá interesa el consumo, que es estable, no el tiempo. Clust. es el tiempo de clustering; Heap, el heap vivo tras cargar el CSV; RSS, el máximo del proceso. CPU es usuario más sistema del proceso completo, incluida la carga del CSV; Núcleos son los núcleos efectivos, CPU dividido por el tiempo de pared total.\par}
\end{table}


Tres observaciones.

\begin{itemize}
  \item \textbf{El pico de memoria es de la carga, no del clustering.} El \emph{heap} tras leer el
        CSV son 349~MB para 136~MB de datos: la diferencia es el crecimiento del arreglo al ir
        agregando filas. Atribuirle ese pico al K-means sería un error de lectura. El RSS máximo
        sube un poco con $P$ por los acumuladores parciales y las pilas de las goroutines, pero la
        memoria la manda el dato. Analíticamente, los datos son $O(n \cdot d)$, los parciales
        $O(C \cdot k \cdot d)$ y las etiquetas $O(n)$.
  \item \textbf{El costo del paralelismo se ve en el tiempo de CPU.} Sube de 5 a 9 segundos mientras
        el tiempo de pared baja. Se paga trabajo extra, sobre todo en coordinar, a cambio de
        terminar antes.
  \item \textbf{Los núcleos efectivos} (CPU dividido por tiempo de pared) llegan a 3 sobre 4 núcleos
        físicos y no se mueven entre $P = 8$ y $P = 16$. Es una comprobación independiente del
        speedup, y coincide: ahí está el techo de la máquina. Que no llegue a 4 se explica porque la
        carga del CSV, que es secuencial, entra en el tiempo de pared total.
\end{itemize}

\subsection{Punto de equilibrio}\label{sec:equilibrio}

El punto de equilibrio es el tamaño de problema a partir del cual la versión concurrente gana a la
secuencial de forma consistente. Se buscó con un barrido fino entre diez mil y cien mil viajes, con
$P = 4$ y el mismo protocolo de \cifra{repeticiones} repeticiones (Tabla~\ref{tab:equilibrio}).

% Generado por tp/scripts/tablas_informe.py desde tp/reports/*.json. No editar a mano.
\begin{table}[H]
\caption{Punto de equilibrio: speedup con $P=4$ según el tamaño del problema}\label{tab:equilibrio}
\footnotesize
\begin{tabular}{@{}rrrll@{}}
\toprule
\textbf{$n$ (viajes)} & \textbf{Secuencial (ms)} & \textbf{Conc. $P=4$ (ms)} & \textbf{Speedup [IC 95\,\%]} & \textbf{¿Gana?} \\
\midrule
1.000 & 0,48 & 0,51 & 0,94$\times$ [0,87; 1,01] & indistinguible \\
10.000 & 4,38 & 4,70 & 0,93$\times$ [0,91; 0,96] & no, pierde \\
15.000 & 6,54 & 6,67 & 0,98$\times$ [0,95; 1,01] & indistinguible \\
20.000 & 8,01 & 7,32 & 1,09$\times$ [1,07; 1,12] & sí \\
30.000 & 12,66 & 7,33 & 1,73$\times$ [1,65; 1,77] & sí \\
50.000 & 20,15 & 8,65 & 2,33$\times$ [2,22; 2,45] & sí \\
70.000 & 29,07 & 9,84 & 2,96$\times$ [2,80; 3,11] & sí \\
100.000 & 41,83 & 14,56 & 2,87$\times$ [2,81; 2,92] & sí \\
\bottomrule
\end{tabular}

{\footnotesize\textit{Nota.} «Gana» cuando el intervalo de confianza queda entero por encima de 1; «pierde» cuando queda entero por debajo. Mismo protocolo que la Tabla~\ref{tab:speedup-fuerte}.\par}
\end{table}


El punto de equilibrio está en \textbf{$n \approx \cifra{n-equilibrio}$ viajes}: el primer tamaño
cuyo intervalo de confianza queda completamente por encima de 1. Por debajo de quince mil viajes la
versión concurrente \emph{pierde} de forma consistente, y el intervalo lo confirma: no es ruido, es
que el costo fijo de armar el pool y repartir bloques supera el trabajo útil de cada iteración, que
con diez mil viajes y $k = \cifra{k}$ son apenas unas decenas de miles de distancias.

Dicho de otro modo: con el dataset del trabajo, \cifra{n}~viajes, unas 140 veces el punto de
equilibrio, la concurrencia está plenamente justificada; no lo estaría para un mes de una flota
chica, y en ese caso la versión secuencial sería la elección correcta por más simple.
